\subsection{Der entscheidende Selektor: Vollständigkeit der elementaren Ströme}

Eine allgemeinere und operativ verständliche Forderung als vorab vorgeschriebener diagonaler Glue lautet:
\begin{equation}\label{eq:currentcompleteness}
 \boxed{\mathcal L=\Z\langle r\in\mathcal L:\|r\|^2=2\rangle .}
\end{equation}
Der vollständige Abschluss soll also von seinen Wurzeln erzeugt werden. Für die hier verwendete Gitter-VOA ist dies gleichbedeutend mit
\[
 V_{\mathcal L}=\langle(V_{\mathcal L})_1\rangle_{\rm VOA}:
\]
alle Felder entstehen durch Produkte, Ableitungen und Linearkombinationen der Ströme vom \emph{konformen Gewicht eins}. Das Gewicht ist eine Eigenschaft der zweidimensionalen Randbeschreibung. Es wird nicht mit dem Spin eines vierdimensionalen Materieteilchens gleichgesetzt.

\textbf{Satz 2.} In der angegebenen gesamten Vergleichsklasse folgt aus dieser Stromvollständigkeit
\begin{equation}\label{eq:globalcurrentselection}
 \boxed{d=4,\quad R_d=A_3\oplus D_5,\quad\mathcal L\cong E_8.}
\end{equation}

\textbf{Beweis ohne Suchgrenze.} Nach Satz 1 genügt $d=4m^2$. Für $m>1$ gilt $d\ge16$. Ein neuer Norm-2-Vektor müsste in einem nichttrivialen Diskriminanten-Coset liegen. Die minimalen Normen auf der $D_{d+1}$-Seite sind
\[
 \mu_D(0)=0,\qquad\mu_D(v)=1,
 \qquad\mu_D(s)=\mu_D(c)=\frac{d+1}{4}>2.
\]
Auf der $A$-Seite hat die Klasse $ka$, $0\le k<d$, die minimale Norm
\[
 \mu_A(k)=\frac{k(d-k)}d.
\]
Da ein isotroper Coset in einem geraden Gitter ausschließlich gerade Normen trägt, kann sein nichttriviales Minimum nicht zwischen null und zwei liegen. Für eine neue Wurzel bleiben deshalb nur folgende Möglichkeiten:
\begin{enumerate}
\item Ein $D$-Spinor-Coset ist schon für sich zu lang.
\item Ein $D$-Vektor-Coset benötigt $\mu_A(k)=1$. Die Gleichung
\[
 k(d-k)=d,\qquad d=k+1+\frac1{k-1}
\]
hat als positive ganzzahlige Lösung mit $0<k<d$ nur $k=2,d=4$.
\item Ein reiner $A$-Coset benötigt $\mu_A(k)=2$. Nun gilt
\[
 k(d-k)=2d,\qquad d=k+2+\frac4{k-2}.
\]
Wegen $k-2\mid4$ ergeben sich nur $(d,k)=(9,3),(8,4),(9,6)$. Keines gehört zu $d=4m^2\ge16$.
\item Ein reiner $D$-Vektor-Coset hat ungerade Normklasse und ist nicht isotrop. Der Null-Coset liefert lediglich die schon vorhandenen Wurzeln von $R_d$.
\end{enumerate}
Für jede mögliche isotrope Glue-Untergruppe und jedes $m>1$ liegen somit alle Norm-2-Wurzeln bereits in $R_d$. Sie erzeugen nur $R_d$, während $[\mathcal L:R_d]=4m>1$. Das widerspricht Gleichung~\eqref{eq:currentcompleteness}. Für $m=1$ ergeben die beiden diagonalen $\Z_4$-Glues E8, das tatsächlich von seinen 240 Wurzeln erzeugt wird. $\square$

Im Gitterfeldbild ist derselbe Beweis anschaulich: Die Ladung eines Feldprodukts ist die Summe der Ladungen seiner Faktoren. Ströme können daher keine Erweiterungsklasse erzeugen, in der überhaupt keine Stromladung liegt. Für die größeren Abschlüsse erscheinen zusätzliche unabhängige Erweiterungssektoren erst bei höherem konformem Gewicht. Dass sie trotzdem holomorph und konsistent sind, macht sie noch nicht zu einer ausschließlich stromerzeugten Naht.

\picturetext{Die größeren Baupläne passen zwar widerspruchsfrei zusammen. Um sie vollständig zu betreiben, braucht man aber zusätzliche elementare Vorgänge, die nicht aus den vorhandenen Strombewegungen entstehen. Beim Viermarkenplan entstehen sämtliche Ergänzungen auf derselben elementaren Ebene. Unter der Forderung, dass diese Ebene vollständig sein soll, bleibt genau dieser Plan.}

\subsection{Warum das an den Originalansatz anschließt -- und was hinzugefügt wird}

Im Original gibt es eine konkrete Vorstufe dieser Forderung: \file{v502\_celestial\_wp5e\_prefactor\_level.py} unterscheidet ganzzahligen Spin von der stärkeren Forderung, dass die Erweiterung auf Strömen mit $h(J)=1$ schließt. Dort wird jedoch bereits die $D_5+A_3$/E8-Architektur verwendet. \file{SEAM.SIMPLECURRENT.GENERATOR.01} verlangt zudem einen tatsächlichen lokalen Erweiterungsgenerator; ein reiner Gitterzeuge genügt seiner analytischen Realisierung nicht.

Die heutige neue Leistung ist die \emph{globale Anwendung} der Stromforderung auf alle $d$, mit vollständigem Ausschluss der unendlichen größeren Familie. Der Standardbegriff von Wurzel- beziehungsweise Stromerzeugung ist keine neue Mathematik. Seine Rolle als gemeinsamer TFPT-Auswahltest wurde hier präzisiert. Die Aussage ist weder als Lean-formalisiert noch als bereits bewiesene Folge von P1/P2 ausgegeben.

Die erste neue Ursprungsforderung ist ausdrücklich Gleichung~\eqref{eq:currentcompleteness}. Reflexionspositivität, Holomorphie, ganzzahliges Gewicht und ein primitiver Windungswert allein beweisen sie nicht. Insbesondere bedeutet ``primitiver Randkernel'' nicht automatisch ``primitive Gittereinbettung''. Die Herkunftsfrage ist jetzt konkret: \emph{Warum erzeugen die lokalen elementaren Nahtoperationen den vollständigen Abschluss, statt nur einen echten Unterraum davon?} Dafür muss man die ursprünglichen lokalen Generatoren und ihre Produkte nachweisen. Weitere Rechnungen an einem bereits gewählten E8-Vakuum würden diese Frage nicht beantworten.

Eine zweite ursprüngliche Route bleibt erhalten: Bei einer flachen Kugel mit ausschließlich $\Z_2$-Kegeldefekten folgt aus Gauss--Bonnet $d\pi=4\pi$, also ebenfalls $d=4$ (\file{QGEO.MARKS.03}). Dort sind Flatness und Defekttyp die Voraussetzungen. Die Übereinstimmung beider Routen ist eine echte gemeinsame Bedingung. Sie beweist noch nicht, dass ihre jeweiligen Voraussetzungen aus einer voraussetzungslosen Naturauswahl folgen.

Die intensive Prüfung der Originale präzisiert dabei die Kausalrichtung. Vier Marken gleicher Kegelordnung sind unter diesen geometrischen Voraussetzungen bereits hergeleitet; sie werden nicht erst aus der späteren E8-Zahl zurückgelesen. Ihre \emph{Lage} bleibt jedoch durch ein Kreuzverhältnis veränderlich. Die Viererfamilie $\{0,1,\lambda,\infty\}$ mit $\lambda\ne0,1$ besitzt stets dieselben Defizite und dieselbe P1-Normierung. Das Quadrat wird erst ausgewählt, wenn die ursprüngliche P2-Clock als konforme Deckwirkung dieser rohen Naht identifiziert wird, \file{QGEO.SYM.01}. Dann ist die freie Viererbahn $\{a,ia,-a,-ia\}$ mit Kreuzverhältnis zwei festgelegt. Die aktuelle Architektur verdichtet den verbleibenden Eingang auf den Symmetrie-Lift-/Twistklassen-Bit, einschließlich der Flagwirkung auf beide Seiten der Marken.

Diese Unterscheidung erhält die vorhandenen physikalischen Herleitungen: $\varphi_0$, die vierfachen $c_3$-Einsetzungen und die Quillen-/BFK-Auswertung benutzen die bereits verbundene Viermarkengeometrie. Ihre rekursive Konsistenz prüft sie stark, liefert aber ohne zusätzlichen Vorwärtssatz nicht die ursprünglich gewählte konforme Deckwirkung. Gleiche Kegelordnung, gleichmäßige modulare Response und gleichmäßige Zustandsgewichte sind drei verschiedene Bedingungen. Die Quellen \file{v216_marks_gauss_bonnet.py}, \file{v181_clock_is_conformal_symmetry.py}, \file{v813_p1_index_kms.py} und die P2-Reduktion in \file{tfpt_1_architecture_e8.tex:1461-1503} unterscheiden sie ausdrücklich.

\subsection{Die stärkere Auswahl braucht nur einen gemeinsamen Strom}\label{sec:onespinorcurrent}

\textbf{Satz 3.} Sei $\mathcal L$ ein positiver gerader unimodularer Abschluss der angegebenen Familie $R_d$. Enthält die zugehörige Gitterfeldalgebra wenigstens ein nichtverschwindendes Feld vom Gewicht eins in einem $D_{d+1}$-Spinorsektor, dann gilt
\begin{equation}\label{eq:onespinorselector}
 \boxed{d=4,\qquad \mathcal L\cong E_8.}
\end{equation}
Weder ein bijektiver Glue noch die Erzeugung der gesamten Feldalgebra durch Ströme wird vorausgesetzt. Auch die konkrete $A$-Diskriminantenklasse muss nicht vorab gewählt werden.

\textbf{Beweis.} Satz 1 liefert $d=4m^2$. Für $m>1$ ist schon der niedrigste mögliche Spinorbeitrag
\[
 h_D(s)=\frac{d+1}{8}\ge\frac{17}{8}>1.
\]
Der $A$-Beitrag und Oszillatoranregungen sind nichtnegativ. Ein spinorgeladenes Gewicht-1-Feld ist daher unmöglich. Es bleibt $m=1$, also $d=4$ und der bereits bestimmte E8-Abschluss. $\square$

Diese Verbesserung verändert den Herkunftstest. Es ist nicht mehr nötig, eine globale Stromvollständigkeit aus P1/P2 herzuleiten. Es genügt, \emph{einen} ursprünglich erzeugten lokalen Strom nachzuweisen, der die Spinorseite und die Markenseite gemeinsam trägt.

Für die elementare Ein-Schritt-Klasse $(s,\omega_1)$ lässt sich die Auswahl sogar ohne Satz 1 direkt schreiben:
\begin{equation}\label{eq:primitivecurrentweight}
 h_d=\frac{d+1}{8}+\frac{d-1}{2d},\qquad
 h_d-1=\frac{(d-4)(d+1)}{8d}.
\end{equation}
Damit gilt $h_d=1$ genau für $d=4$. Schon ganzzahliges $h_d$ verlangt $d\mid4$; $d=2$ hat jedoch $h_d=5/8$. Diese Formel betrifft $\omega_1$ beziehungsweise $\omega_{d-1}$, nicht jeden primitiven Generator von $\Z_d$. Für eine allgemeine Klasse $k$ ist $h_A(k)=k(d-k)/(2d)$.

\subsection{Was Skalenfreiheit tatsächlich erzwingen würde}\label{sec:seamwardsource}

Die Originalquelle liefert $[u_\Sigma]=1\in K^1(S^1)$ und Spektralfluss eins. Das ist ein topologischer Transportwert. Der normierte Seam-Modus und der im Original eingeführte Projektor $|\psi_\Sigma\rangle\langle\psi_\Sigma|$ sind noch kein lokaler chiraler Strom. Auch der Calderón-Projektor bestimmt zunächst Cauchy-Daten und Polarisation.

Eine konkrete hinreichende Herkunftsidentität wäre
\begin{equation}\label{eq:seamcontouridentity}
 Q_\Sigma=\frac1{2\pi i}\oint\Psi(z)\,dz\ne0,
 \qquad [L_0,Q_\Sigma]=0,
\end{equation}
wobei $\Psi$ ein lokales chirales Primärfeld im Spinorsektor ist und der Konturtransport die ursprüngliche Nahtoperation implementiert. Vorausgesetzt wird zudem, dass dieser nichtverschwindende Sektorintertwiner auf dem ausgewählten Raum entlang der geschlossenen Kontur eindeutig einwertig definiert ist: Die Integration besitzt keinen zusätzlichen Monodromie-Randterm. Der Konturwert bleibt bei Radiusänderung derselbe; die ausdrücklich geforderte interne Dilatationsinvarianz bleibt davon als eigene Bedingung zu unterscheiden. Die Primärfeldidentität
\[
 [L_0,\Psi(z)]=(z\partial_z+h)\Psi(z)
\]
ergibt durch Integration
\[
 [L_0,Q_\Sigma]=(h-1)Q_\Sigma.
\]
Folglich ist $h=1$, und Satz 3 wählt $d=4$. Hier wird eine \emph{ungewichtete} interne Konturoperation verlangt. Der gewichtete Nullmodus $\oint z^{h-1}\Psi(z)\,dz$ kommutiert für jedes ganzzahlige $h$ mit $L_0$; beispielsweise ist $L_0=\oint zT(z)\,dz$ dimensionslos, obwohl $T$ Gewicht zwei besitzt. Eine mitgeführte Radiusnormierung oder eine laufende Kopplung würde den Schluss ebenfalls verändern. Für eine tatsächlich lineare eindimensionale Kopplung $g\oint d\ell\,\Psi$ gilt dieselbe Bedingung als $[g]=1-h$; eine zweidimensionale Wechselwirkung hat eine andere Dimensionszählung.

Der Universalraum muss deshalb unterscheiden, ob seine Rekursion \emph{dieselbe interne Operation} bei allen Radien erhält oder lediglich eine mit ihrer Skala transformierte Operation. Antworttreue allein setzt diese beiden Fälle nicht gleich. Ein endlicher Vierteldrehungsdefekt kann zudem konforme Symmetrien respektieren, ohne der Nullmodus eines kontinuierlichen Stroms zu sein.

Winding und Glueklasse bestimmen das absolute Gewicht nicht. Im selben ursprünglichen E8-Coset haben $\lambda$ und $\lambda+r$, mit $r=(1,-1,0,0,0;0,0,0)$, dieselbe Ordnung-4-Klasse, aber $h_\lambda=1$ und $h_{\lambda+r}=2$. Die vorhandenen Alpha- und Flavorantworten müssen daher an den tatsächlich erzeugten lokalen Träger angeschlossen werden. Die Identität $4\cdot2\pi\cdot c_3=1$ ersetzt weder diesen Stromnachweis noch den ursprünglichen Determinantenlinien-Response.

\subsection{Der konstruktive Anschluss: Zwei Vierteltwists ergeben das E8-Feld}\label{sec:doubletwistbridge}

Die bedingte Achtkanal-Fermionisierung enthält einen viel direkteren Kandidaten für diesen Strom. Sechzehn reelle CAR-Fermionen ergeben acht unabhängige komplexe Kanäle. In der gewählten $D_5\oplus D_3$-Zielmarkierung werden sie als $10+6$, also fünf plus drei orthogonale Bosonisierungskoordinaten, organisiert. Die Kopienzahl allein leitet diese Blockwirkung nicht her. Die vier üblichen $A_3$-Koordinaten sind dagegen durch ihre Summe null eingeschränkt. Sie dürfen nicht als vier zusätzliche unabhängige Fermionenkanäle gelesen werden. Die Isometrie, der Halb-Twist und sein ganzzahliger Ladungs-Carry wurden bereits im Vertrag \file{half-twist-grade-carry} vom 8.~September ausgearbeitet. Sie sind keine erst heute gefundenen Quellenresultate. Neu in der Gesamtsynthese sind ihre Verknüpfung mit Satz 3 und die explizite OPE-Form der gesuchten Herkunftsidentität.

Diese Voraussetzung ist auch mikroskopisch zu beachten: Der MMST-Text \file{mmst\_charged\_scaling\_limit\_en.tex} setzt sechzehn Majorana-Kopien an. Der unveränderte einzelne QWZ-Kragen besitzt dagegen einen komplexen Kanal je Rand; $N_y=8$ bezeichnet seine Querbreite und liefert keine acht Spezies. Das wurde im vorhandenen Vertrag \file{source-rg-clock-bridge-20260920} analytisch geprüft. Der in Abschnitt~\ref{sec:nativeclosure} konstruierte ursprüngliche globale Träger $V_{16}$ besitzt die passende Darstellung $10+6$; dass gerade der mikroskopische Nahtoperator auf diesem gesamten Träger wirkt, ist damit noch nicht bewiesen. Ein bloßes Hinzufügen von Kopien würde diese Herkunftsidentifikation nicht liefern.

Die explizite Isometrie $A_3\to D_3$ ist
\begin{equation}\label{eq:a3d3twistisometry}
 T=\frac12
 \begin{pmatrix}
 1&1&-1&-1\\
 1&-1&1&-1\\
 1&-1&-1&1
 \end{pmatrix},\qquad
 TT^{\mathsf T}=I_3,\qquad
 T^{\mathsf T}T=I_4-\frac14\mathbf1\mathbf1^{\mathsf T}.
\end{equation}
Sie bildet die zwölf $A_3$-Wurzeln genau auf die zwölf $D_3$-Wurzeln ab, nicht lediglich auf einen gleichdimensionalen Raum. Insbesondere gilt
\[
 T\left(\frac34,-\frac14,-\frac14,-\frac14\right)
   =\left(\frac12,\frac12,\frac12\right).
\]
Der ursprüngliche gemeinsame Vektor $\lambda=(\omega_s,\omega_f)$ hat in den \emph{acht unabhängigen} Kanälen deshalb die Form
\begin{equation}\label{eq:twistdoublecharge}
 \lambda_{8}=\left(\frac12\right)^8,
 \qquad q_{1/4}=\left(\frac14\right)^8=\frac12\lambda_8,
 \qquad h(q_{1/4})=\frac14,\quad h(\lambda_8)=1.
\end{equation}
Der einfache Vierteltwist ist weiterhin nicht das Gewicht-1-Feld. Er besitzt aber exakt die halbe Ladung des gesuchten Feldes. Im erklärten freien chiralen Bosonisierungsrahmen, mit $\phi_i(z)\phi_j(w)\sim-\delta_{ij}\log(z-w)$, folgt für die Twist-Intertwiner $V_q=:\!e^{iq\cdot\phi}\!:$
\begin{equation}\label{eq:doubletwistope}
 V_q(z)V_q(w)
 =c(q,q)(z-w)^{1/2}
 \bigl[V_{\lambda_8}(w)+\text{Terme höherer Ordnung}\bigr],
 \qquad c(q,q)\ne0.
\end{equation}
Das ist eine konkrete OPE-Konstruktion: Nach Abziehen des bekannten Abstandsfaktors ist der führende Fusionskoeffizient gerade das E8-Feld. Eine Wahl des Intertwiner-Kokzyklus bestimmt die Phase $c(q,q)$. Die Gleichung behauptet keinen unbelegten starken Operatorgrenzwert für die mikroskopischen QWZ-Strings.

Die Lokalität und der vollständige algebraische Anschluss lassen sich in der vorhandenen E8-Ziel- und Chevalley-Algebra prüfen. Für jede Wurzel $r$ von $D_5\oplus D_3$ ist $\lambda_8\cdot r\in\Z$; das Spinorfeld ist daher relativ zu diesen geraden Strömen lokal. Die Wurzeln dieses Basissektors und ein solches Feld samt Adjungiertem erzeugen unter den tatsächlichen Chevalley-Klammern alle 240 E8-Wurzeln. Die acht Cartanrichtungen ergeben zusammen 248 Ströme. Die Zerlegung ist
\[
 248=(45,1)\oplus(1,15)\oplus(10,6)
      \oplus(16,4)\oplus(\overline{16},\overline4).
\]
Dabei kommen 52 Wurzeln aus der Basis, 60 aus dem gemischten Vektorsektor und 128 aus einer gemeinsamen Spinorchiralität. Eine einzelne ursprüngliche CAR-Fermionenladung $e_i$ hat hingegen $\lambda_8\cdot e_i=1/2$. Das neue Feld ist nicht lokal gegenüber der ganzen unprojizierten CAR-Feldalgebra. Auch beide entgegengesetzten Spinorchiralitäten dürfen nicht gleichzeitig als lokale bosonische Erweiterung aufgenommen werden: ein entsprechendes Gegenpaar hat inneres Produkt $3/2$. Die Paritäts-/GSO-Auswahl ist ein wirklicher Teil der Konstruktion.

Der bekannte Zusammenhang zwischen 16 chiralen Majorana-Fermionen und dem E8-Level-1-Modell wird durch die vollständigen NS-/R-Sektoren, nicht durch die Zahl 16 allein, beschrieben~\cite[Abschnitt 5.1, Gleichungen (5.3)--(5.6)]{boylesmith}. Die vorstehende Isometrie und OPE geben dazu das konkrete Wörterbuch für den bereits vorhandenen TFPT-Generator. Eine neue Raumzeit oder neue Materiefamilie wird dafür nicht eingeführt.

\picturetext{Ein Viertelschritt ist noch kein vollständiger Strom. Zwei Viertelschritt-Felder besitzen zusammen genau die Ladung des fehlenden Verbindungsstücks. Mit diesem Stück erreichen die vorhandenen Ströme alle Teile des E8-Gerüsts. Man muss dabei das Feldprodukt mitsamt Abstand, Phase und Sektor behandeln; zweimal dieselbe Clocktaste zu drücken ist mathematisch eine andere Operation.}

Gerade diese Unterscheidung verhindert eine falsche Gleichsetzung: Im zusätzlichen vierstufigen Fluxregister von \file{v1033\_charged\_disorder.py} trägt das Quadrat des Shifts die Registerladung zwei. Der Vektor $\lambda$ hat gegenüber $D_5\oplus A_3$ dagegen Glueklasse eins modulo vier. Sein Monodromiezeichen auf den acht elementaren Fermionen ist $-1$ und hat Ordnung zwei. Registerladung, Fermionenmonodromie und interne Gluegraduierung sind drei verschiedene Wirkungen. Ein geänderter Basisname identifiziert sie nicht.

Die bereits vorhandene Ladungsdarstellung macht dies auf allen Gitterzuständen präzise. Mit $s=\lambda_8$ gilt
\[
 L_0=D_5\oplus D_3\subset D_8\subset L=D_8\cup(D_8+s),
 \qquad L/L_0\cong\Z_4,\quad L/D_8\cong\Z_2.
\]
Der interne Grad ist $k(p)=2\sum_{i=1}^5p_i\bmod4$. Ein geladener Nullmodenfaktor wirkt dagegen auf dem unendlichen Raum $\ell^2(L)$ als
\[
 U_s\ket p=\varepsilon(s,p)\ket{p+s},\qquad
 CU_sC^{-1}=iU_s,\quad C\ket p=i^{k(p)}\ket p.
\]
Im im Vertrag festgehaltenen Kokzyklus ist $U_s^2=-U_{2s}$ und $U_s^4=U_{4s}\ne I$. Die skalare Phase der ersten Gleichung hängt vom Kokzyklus ab; der tatsächliche ganzzahlige Ladungsübertrag der zweiten nicht. Nach vier Schritten ist die Gradmarke zurück am Anfang, die Ladung aber um $4s$ verschoben. Der endliche Registerumlauf allein würde diese Wirkung verlieren. Die Oszillatorfaktoren des vollständigen Vertexfeldes bleiben zusätzlich nötig.

\picturetext{Ein Zeiger zeigt nach vier Viertelschritten wieder in dieselbe Richtung. Das Fahrzeug ist trotzdem weitergefahren. Der Universalraum muss Zeigerstellung und zurückgelegten Weg erhalten, sobald spätere Antworten beide unterscheiden. Ein Zuse- oder Hypergraphmodell, das nur die vier Zeigerstellungen speichert, würde genau den Ladungs-Carry und damit das vollständige Nahtfeld verlieren.}

\textbf{Damit ist der Anschluss konstruktiv, aber noch bedingt.} Das fehlende mikroskopische Feld ist nicht länger nur abstrakt mit $h=1$ zu fordern: Es ist als führender Fusionskoeffizient zweier passender Twist-Intertwiner benannt. Für die Originalquelle sind die support-erhaltende Realisierung dieses Fusionskoeffizienten, sein nichtverschwindender Grenzwert, die gemeinsame Kokzyklus-/Chiralitätswahl und der Abstieg aus dem zusätzlichen Fluxregister zu beweisen. Der vorhandene einfache Vierteltwist-Audit wird dadurch nicht widerlegt, und der offene MMST-/FE-GEN-Vertrag wird nicht geschlossen. Die A3/D3-Isometrie ist speziell an der schon ausgewählten Viermarkenquelle verfügbar; sie darf nicht rückwärts als voraussetzungslose Auswahl von vier Marken ausgegeben werden. Satz 3 beschreibt, weshalb eine unabhängig hergeleitete solche Spinorstrom-Operation die globale Auswahl tatsächlich schließen würde.

\subsection{Der Mikrounterbau ist weiter geschlossen als zuvor ausgewiesen}\label{sec:microscopicpairbridge}

Die erneute Quellenprüfung korrigiert eine zu schwache Darstellung des vorhandenen Bestands. Der Vertrag \file{microscopic-neutral-limit} beweist für den unveränderten QWZ-Kragen mit Breite acht, Masse eins, Holonomiesektor $r=1$, tatsächlichem gefülltem negativem Spektrum und Filterbreite $\delta=4N^{-3/4}$ bereits
\[
 I_N=O(N^{-1/4})=o(N^{-3/16}),\qquad
 Z_t^2(1+2\sqrt2A_{\rm ref})I_N=O(N^{-1/16}),
 \qquad t=N\delta.
\]
Dies ist der analytische Vergleich der \emph{vollständigen} bewegten Geschichte, einschließlich der Diagonalrahmen und des unteren Rands. Der ältere Vertrag \file{history-reference-transport} formulierte diese Rate noch als offen; der nachfolgende Beweis schließt sie. Der Vertrag \file{microscopic-fourpoint-limit} erweitert die Antwortkontrolle mit einem logarithmischen Hilfsfenster und erhält geordnete Vierpunktantworten einschließlich ihrer Kollisionsgrenzen. Die Quelldynamik wird dabei nicht geändert.

Auch das ganzzahlig geladene lokale Feld ist bereits konstruiert: \file{microscopic-charged-car-limit} benutzt die echten Erzeugungs- und Vernichtungsoperatoren auf der oberen Randzeile, ihre wirklichen Adjungierten, das gefüllte Quellvakuum und dieselbe Zeit. Es beweist Konvergenz aller festen endlichen Feldwörter und Zeiten sowie einen gemeinsamen positiven Energiekern. Das ist mehr als ein neutraler Zahlenvergleich. Es liefert einen komplexen chiralen CAR-Kanal; es identifiziert weder P1s rohen Kernel mit diesem gewählten QWZ-Modell noch diesen einen Kanal mit dem gesamten $V_{16}$.

\textbf{Eine neue Folgerung verbindet diese Ergebnisse mit der gesuchten Gewicht-1-Paarantwort.} Vorausgesetzt seien genau die bereits deklarierten acht identischen, gleich orientierten QWZ-Kopien mit Produktvakuum. Keine Kopie wird aus der transversalen Breite abgeleitet. Für eine Kopie sei
\[
 C_N^\circ(a,b)=
 e^{\,i\Tr(P_NF_{a,N})-i\Tr(P_NF_{b,N})}
 \det_{\operatorname{ran}P_N}
 \left(P_Ne^{-iF_{a,N}}e^{iF_{b,N}}P_N\right)
\]
das vorhandene normalgeordnete Zweipunktwort auf dem \emph{vollen} besetzten Quellraum. Die Phasenkorrektur ist eine feste Endpunkt-Koboundary. Sie erhält die Grampositivität und ist keine angepasste Phase. Beide rohen Vakuumüberlappungen, Quelle und Stromreferenz, haben Betrag höchstens eins.

Der bereits bewiesene Vergleich gilt für zwei Beine ebenso wie für vier. Im unveränderten logarithmischen Hilfsfenster
\[
 M=\Theta(N^{1/4}\log N),\qquad
 I_{\rm lin}+I_{\sharp,\rm ref}=O(N^{-1}(\log N)^3),
 \qquad A_{\rm ref}=O(N^{1/8})
\]
liefert er vor der Endpunktnormierung
\begin{equation}\label{eq:onecopyrawpairrate}
 \sup_{\rm grid}|C_N^\circ-C_{{\rm cur},t}^\circ|
 =O(N^{-7/8}(\log N)^3).
\end{equation}
Der separate endliche Stromrest ist kleiner. Es handelt sich nicht um eine neue Schätzung, die Diagonalrahmen oder Bulkmoden wegwirft: Diese Rate wird aus den vorhandenen vollständigen Geschichtsabschätzungen gewonnen.

Setze $w_k(t)=e^{-(2\pi k/t)^2}$,
$H_t=\sum_{k\ge1}w_k(t)/k$ und $Z_t=e^{H_t/8}$. Das unabhängige Kopienprodukt und die teleskopische Identität für achte Potenzen ergeben
\[
 |(C_N^\circ)^8-(C_{{\rm cur},t}^\circ)^8|
 \le8|C_N^\circ-C_{{\rm cur},t}^\circ|,
 \qquad Z_t^{16}=e^{2H_t}=O(t^2)=O(N^{1/2}).
\]
Folglich gilt für die echte neutrale Paarantwort der acht Kopien
\begin{equation}\label{eq:eightcopyspinorpairrate}
 \boxed{\sup_{\rm grid}
 \left|Z_t^{16}(C_N^\circ(a,b))^8-\mathcal K_t(b-a)\right|
 =O(N^{-3/8}(\log N)^3),}
 \qquad
 \mathcal K_t(s)=
 \exp\!\left(2\sum_{k\ge1}\frac{w_k(t)}k e^{2\pi i k s}\right).
\end{equation}
Die achte Potenz entsteht hier aus acht unabhängigen Quelldeterminanten, nicht aus dem Potenzieren eines bereits genommenen singulären $L^1$-Grenzwerts.

Die Kollisionsgrenze lässt sich direkt und ohne diese unzulässige Vertauschung führen. Die analytische Potenzreihe $\mathcal K_t(z)=\sum_{n\ge0}a_n(t)z^n$ hat $a_0=1$ und
\[
 n a_n(t)=2\sum_{k=1}^n w_k(t)a_{n-k}(t).
\]
Weil $0\le w_k\le1$, folgt induktiv
\[
 0\le a_n(t)\le n+1,\qquad a_n(t)\longrightarrow n+1.
\]
Die Majorante $(n+1)$ ist gegen die rasch fallenden Fourierkoeffizienten jedes glatten Testfeldes summierbar. Daher konvergiert die Stromreferenz als Distribution auf dem Kreis. Zudem ist
\[
 \sup|\mathcal K_t|\le e^{2H_t}=O(t^2),\qquad
 \sup|\mathcal K_t'|=O(t^3).
\]
Für den Vergleich auf dem ganzen Endpunktquadrat wird der mikroskopische Kernel ausdrücklich stückweise konstant fortgesetzt:
\[
 \widehat a_N=\frac{\lfloor Na\rfloor}{N},\qquad
 \widehat b_N=\frac{\lfloor Nb\rfloor}{N},\qquad
 \mathcal K_N^{\rm pc}(a,b)
 =Z_t^{16}\bigl[C_N^\circ(\widehat a_N,\widehat b_N)\bigr]^8.
\]
Die Rundung beider Endpunkte auf das Quellgitter kostet höchstens $O(t^3/N)=O(N^{-1/4})$. Zusammen mit Gleichung~\eqref{eq:eightcopyspinorpairrate} folgt somit einschließlich der Kollisionsdiagonale als Teil der distributionellen Aussage
\begin{equation}\label{eq:microscopicweightonepair}
 \boxed{\mathcal K_N^{\rm pc}
 \ \longrightarrow\
 (1-e^{2\pi i(b-a)-0})^{-2}
 \quad\text{in }\mathcal D'(\mathbb T^2).}
\end{equation}
Der radiale Randwert ist festgehalten. Positivität bleibt sowohl im mikroskopischen Tensor-Gramkernel als auch durch die nichtnegativen Fourierkoeffizienten erhalten.

\textbf{Das liefert die Gewicht-1-Kernform, noch kein lokales geladenes Gewicht-1-Feld.} Der Kern stimmt mit der Zweipunktform des konstruierten Spinorstroms überein. Ein Zweipunktkern allein bestimmt jedoch weder Nullmodenladung, Kokzyklus, gemeinsame Spinorchiralität, lokale Austauschprodukte noch die Wirkung aller gemischten Ströme. Die $N$-abhängigen Dichteexponentiale sind auf dem ganzen Zylinder insgesamt ladungsneutral. Nach den vorhandenen Twist-Transporten kann ein oberes Teilchen mit einem unteren Loch kompensiert werden. Das schließt einen geladenen Sektor der oberen Randalgebra nicht grundsätzlich aus; der lokalisierte Ein-Rand- und Sektoranschluss muss aber wirklich bewiesen werden. Eine bloße neue Ladungsnummer oder ein angehängtes Fluxregister wäre kein solcher Anschluss.

\picturetext{Das diskrete Instrument liefert bereits die richtige gemeinsame Klangform des gesuchten Verbindungsstücks. Wir müssen nicht erneut fragen, ob überhaupt ein kontinuierlicher Klang entsteht. Zu beweisen ist, welche konkrete lokale Betätigung diesen Klang mit seiner Ladung erzeugt und weshalb die ursprüngliche Geometrie genau die acht beteiligten Kanäle trägt. Erst dann gehört das Stück nachweislich zum ursprünglichen Instrument.}

\subsection{Wie die vorhandenen physikalischen Herleitungen zusammen zurückwirken}

Unter dem neuen Selektor wird das diskrete Gerüst auf einmal festgelegt:
\[
 d=4\ \Longrightarrow\
 (\rank A_3,\rank D_5)=(3,5),\quad
 \dim\Lambda^{\rm even}\C^5=16,\quad
 A_3+D_5\longrightarrow E_8.
\]
Das \emph{zusätzliche originale Auslesewörterbuch} liest die drei Zyklusrichtungen als Familien und den Halbspinor als eine chirale Materiegeneration; damit folgen $N_{\rm fam}=3$ und die $48$ Weylkomponenten. Der globale Gitterbeweis allein macht aus einer Zyklusrichtung kein physisches Teilchen und beweist diese Ausleseidentifikation nicht.
Bei der im Original begründeten positiven ganzzahligen Ankertypisierung folgt aus drei Gewichten mit Summe vier der Anker $(1,1,2)$. Seine elementaren symmetrischen Polynome sind $(4,5,2)$. Der E8-Coxeterwert ist $30=2\cdot3\cdot5$; die Viertelmarkierung wirkt als Galois-Zahnrad, nicht als unzulässige Untergruppe $\Z_4\subset\Z_{30}$.

Die vorhandene Trägerlesart $\C^5=\C^3\oplus\C^2$ und
\[
 Y=-\frac13N_3+\frac12N_2
\]
liefern auf demselben geraden Cliffordträger
\[
 \Tr_{48}Y^2=10,\qquad b_1=\frac{41}{10},\qquad
 \Tr_{\C^5}Y^2=\frac56.
\]
Mit der ursprünglichen P1-Randnormierung folgen die bereits dokumentierten Seed-Beziehungen
\[
 c_3=\frac1{8\pi},\qquad
 \phi_0=\frac43c_3+48c_3^4,
 \qquad\lambda_C=\sqrt{\phi_0(1-\phi_0)}.
\]
Flavorfiltration, sechsfacher Transport, massenwirksame Wortlängen und ladungsgewichtete Alpha-Antwort sind damit eng eingeschränkt. Die neue globale Auswahl erzeugt diese bekannten Herleitungen nicht ein zweites Mal; sie untersucht, warum gerade deren gemeinsames diskretes Gerüst auftreten könnte.

\begin{tabularx}{\linewidth}{lY}
\toprule
Gemeinsamer Anschluss & Konsequenz für eine vollständige Interpretation\\\midrule
Naht, Double Cover, vier Marken & Ein lokaler spinorgeladener Gewicht-1-Strom genügt zur Auswahl. Sein Nahtursprung und seine Sektoren müssen nachgewiesen werden; vier Marken sind keine vier Raumzeitkoordinaten.\\
Träger, Ladungen, Flavor & Familien- und Ladungsauslesungen müssen dieselben Markierungen und die tatsächliche Orientierung respektieren; Ränge und Materiesorten bleiben getrennt.\\
P1, Alpha, geometrische Skalen & Der Gitterindex vier ist mit der P1-Identität $4\cdot2\pi\cdot c_3=1$ exakt kompatibel. Der reale normierte Antwortoperator und sein gemeinsamer Indexanschluss müssen nachgewiesen werden.\\
Neutrinos und CP & Der vorhandene Majoranapaarraum ist erlaubt. Seine Phase und Familienorientierung müssen aus dem Prozess folgen. Der führende Seesaw liefert $\theta_{13}=0$; der separate Readout $\phi_0e^{-5/6}$ muss im vollständigen Mischungsoperator wiedergefunden werden.\\
Higgs und Gravitation & Schwaches Higgsdoublett, neutraler Scalaron und $B-L=2$-Paarsektor sind unterschiedliche Rollen. Eine gemeinsame Quelle erlaubt ihre Kopplung, macht sie vor der Brechung aber nicht zu demselben Feld.\\
Universalraum, Zuse, Hypergraph & Die Darstellung muss die ganzen zusammengesetzten Antworten erhalten. Erhaltene Inzidenz oder Einzeitdaten allein beweisen weder die gewählte Regel noch die räumliche Realisierung.\\\bottomrule
\end{tabularx}

Die Neutrino- und Kopplungsbedingungen sind Rückbedingungen an die \emph{ganze} Interpretation. Sie bestimmen nicht vorab den einzigen nächsten Forschungszweig. Besonders darf die erfolgreiche abelsche Alpha-Gleichung nicht durch bloßen Austausch des Betakoeffizienten auf die nichtabelschen Kopplungen übertragen werden; dieser naive Weg ist im Originalvertrag \file{GAUGE.DETLINE.FIXPOINT.01} ausgeschlossen. Die absolute Einheitenskala bleibt die explizite metrologische Kalibrierung des Originals.

\subsection{Was dies für den Universalraum und die vollständige Lösung bedeutet}

Der Universalraum fordert eine vollständige, zusammensetzbare Darstellung der erlaubten Antworten. Für zwei Geschichten $h,h'$ ist die sichere Gleichsetzung
\[
 h\sim h'\quad\Longleftrightarrow\quad
 p(r\mid\mathcal C[h])=p(r\mid\mathcal C[h'])
 \quad\text{für alle zulässigen Testkontexte $\mathcal C$ und Records $r$}.
\]
Die Testkontexte müssen auch spätere Operationen und zusammengesetzte Umgebungen erfassen. Besitzt die Prozesstheorie eine Parallelkomposition, müssen alle ihre zulässigen parallelen Registerkontexte ebenfalls enthalten sein. Bei Abschluss der Testklasse unter Zusammensetzung wird diese Relation zur entsprechenden Kongruenz. Ein Vergleich bloßer momentaner Zustände genügt nicht. Dieser Quotient komprimiert ein gegebenes vollständiges Gesetz; er wählt das Gesetz nicht aus.

Damit stehen zwei unterschiedliche Vollständigkeiten nebeneinander. Die Stromforderung verlangt, dass die elementaren lokalen Generatoren die ganze algebraische Quelle erzeugen. Der Universalraum verlangt, dass ihre Darstellung keine später unterscheidbare Antwort verliert. Keine der beiden Forderungen ersetzt die andere. Zusammen formulieren sie einen präzisen Kandidaten für ein sparsames Ursprungsprinzip: \emph{vollständige Erzeugung und vollständige Wirkungserhaltung desselben lokalen Prozesses}. Seine Realisierung darf mehrere effektive Räume besitzen und muss die Originalauslesungen über nachgewiesene Abbildungen verbinden.

\begin{center}
\begin{tikzpicture}[node distance=8mm]
\node[box,text width=13cm] (o) {\textbf{Ursprüngliche markierte Naht}\\Zyklenseite und Funktions-/Cliffordseite desselben Objekts};
\node[box,below=of o,text width=13cm] (c) {\textbf{Globale algebraische Auswahl}\\Gerader unimodularer Abschluss: $d=4m^2$\\Ein spinorgeladener Gewicht-1-Strom: ausschließlich $d=4$, E8};
\node[box,below=of c,text width=13cm] (p) {\textbf{Vorhandene gekoppelte Auslesungen}\\Familien, Ladungen, Clocks, Flavor, Alpha und geometrischer Randkanal};
\node[box,below=of p,text width=13cm] (u) {\textbf{Vollständiger physischer Prozess}\\Generatoren, Zustand/Phase, lokale Realisierung und tatsächliche Antworten\\Universalraum erhält ihre zusammengesetzten Wirkungen};
\draw[dasharr] (o)--node[right,font=\scriptsize]{Strom aus der Naht herzuleiten}(c);
\draw[arr] (c)--node[right,font=\scriptsize]{originales Auslesewörterbuch}(p);
\draw[dasharr] (p)--node[right,font=\scriptsize]{physische Identifikation}(u);
\end{tikzpicture}
\end{center}

\textbf{Der Stand ist damit konkreter, aber keine bewiesene Gesamtlösung.} Wir besitzen einen allgemeinen bedingten Ursprungsauswahlsatz, einen engeren Test mit nur einem gemeinsamen Strom und die bereits ausgearbeiteten Zielauslesungen. Die Herkunftsidentität betrifft jetzt einen bestimmten Fusionskoeffizienten aus den bedingt angenommenen achtkanaligen Twist-Intertwinern samt Ladungs-Carry. Die ursprüngliche Nahtoperation und P1-Antwort müssen an ihn angeschlossen werden. Sein tatsächlicher Operatoranschluss muss auch die bekannte Trägerdarstellung $V_{16}$ erhalten. Die Massen-, Flavor-, Eich- und Gravitationsauslesungen sind gemeinsame Rückbedingungen an diesen Prozess. Ein passender Zahlenwert oder eine zusätzliche Clock-Matrix könnte diese Herkunftsidentifikation nicht ersetzen.

\section{Das Gesamtbild: Eine Quelle und mehrere Ansichten derselben Wirklichkeit}\label{sec:bigpicture}

\key{\textbf{Das verbindende Bild ist ein Instrument mit einer einzigen inneren Mechanik.} Topologie legt fest, welche Schleifen und Markierungen unterscheidbar sind. Algebra legt fest, wie Operationen zusammenwirken. Der Zustand bestimmt ihre gemeinsamen Antworten. Rekursion beschreibt, wie dieselben Antworten auf einer anderen Auflösung repräsentiert werden. Physik verlangt zusätzlich, dass dieses Instrument tatsächlich unsere räumlichen, zeitlichen und messbaren Prozesse beschreibt.}

Die erneute Prüfung korrigiert die bisherige Gewichtung dieses Dokuments. Die ursprünglichen Paper haben den Weg von den Markierungen zu physikalischen Größen bereits weit ausgearbeitet. Die jüngere Sitzung hat dagegen vor allem einen gemeinsamen Feld- und Prozesskandidaten untersucht. Das Gesamtbild muss beide Richtungen verbinden: vom Ursprung zu den bekannten Auslesungen und von diesen Auslesungen zurück zu den erlaubten Ursprüngen.

\begin{center}
\begin{tikzpicture}[node distance=7mm]
\node[box,text width=13.5cm] (p) {\textbf{Ursprung und Markierungen}\\P1, P2, Double Cover, $\mu_4$, Anker $(1,1,2)$};
\node[box,below=of p,text width=13.5cm] (l) {\textbf{Gemeinsamer Träger}\\markierte $D_5+A_3$-/E8-Struktur, drei Familien, Ladungen};
\node[box,below=9mm of l,xshift=-3.6cm,text width=6.0cm] (a) {\textbf{Vorhandener Flavorweg}\\Familienmonodromie, $R,Q,\Sigma$, $K,L$\\Quellenmassen, Yukawa, CKM/PMNS};
\node[box,below=9mm of l,xshift=3.6cm,text width=6.0cm] (b) {\textbf{Vorhandener Nahtweg}\\$c_3$, $48$ Zustände, gewichtetes Budget $41$\\$\phi_0$, $\alpha$, Skalenrelationen};
\node[box,below=29mm of a,xshift=3.6cm,text width=13.5cm] (o) {\textbf{Bereits berechnete gemeinsame Bedingungen}\\gleicher Seed, gleiche Markierung, gleiche Normierung; exakte und bedingte Ergebnisse getrennt};
\node[box,below=11mm of o,text width=13.5cm] (s) {\textbf{Gesuchte vollständige gemeinsame Quelle}\\derselbe Operator, Zustand und Prozess müssen beide Wege realisieren\\Die neue E8-Level-1-Quelle ist daran zu prüfen};
\node[box,below=of s,text width=13.5cm] (r) {\textbf{Universalraum und physische Realisierung}\\antworttreue Geschichten, lokale Dynamik, RG-/Messabbildung und Raumzeit};
\draw[arr] (p)--(l);\draw[arr] (l)--(a);\draw[arr] (l)--(b);
\draw[arr] (a)--(o);\draw[arr] (b)--(o);
\draw[dasharr] (o)--node[right,font=\scriptsize]{gemeinsame Auswahl}(s);
\draw[dasharr] (s)--(r);
\end{tikzpicture}
\end{center}

Durchgezogene Pfeile bezeichnen hier vorhandene Herleitungs- und Rechenwege mit den im Original ausgewiesenen Voraussetzungen. Sie erklären nicht jeden Endwert zu einem voraussetzungslosen Satz. Gestrichelte Pfeile bezeichnen die noch nachzuweisende gemeinsame Auswahl beziehungsweise physische Realisierung. Insbesondere liegen Massen und $\alpha$ nicht erst hinter einer vollständig neuen Raumzeittheorie auf einer leeren Aufgabenliste.

\key{\textbf{Die starke Verbindung lautet:} Derselbe Bauplan liefert unterschiedliche, miteinander gekoppelte Antworten. Das ungewichtete Zählen der Materiezustände ergibt $48$, die elektromagnetische Gewichtung zusammen mit dem Higgs ergibt $41$, und der Flavortransport verbraucht genau $40$ Einheiten plus den Higgsbeitrag. Der Nahtparameter bestimmt zugleich Massenverhältnisse und Mischungen. Ein gemeinsamer Ursprung muss diese vollständigen Operatorbeziehungen erhalten.}

\picturetext{Die ursprüngliche Theorie enthält bereits Zahnräder und berechnete Übersetzungsverhältnisse. Die neue Quellenkonstruktion soll zeigen, wie diese Zahnräder auf derselben Welle sitzen und tatsächlich zusammen laufen. Man darf weder die bekannten Übersetzungen vergessen noch aus mehreren einzeln laufenden Rädern schließen, dass sie schon miteinander verbunden sind.}

Der neue Quellkandidat beantwortet ganze Folgen von Feldeinsetzungen. Seine Algebra, sein Zustand und seine konforme Zeit sind unter der erklärten Quantisierung gemeinsam bestimmt. Das liefert einen konkreten Kandidaten für den Universalraum. Ob seine markierten Projektionen tatsächlich die ursprüngliche Flavorfiltration, den elektromagnetischen Determinantenkanal und die physische Dynamik liefern, ist eine stärkere Frage. Sie muss an den schon vorhandenen Antworten entschieden werden.

Eine eindeutige Darstellung einer gewählten Quelle ist noch keine eindeutige Wahl der Quelle. Die ursprünglichen physikalischen Ergebnisse sind dabei zusätzliche Auswahlbedingungen. Sie werden in Abschnitt~\ref{sec:physicalconstraints} ausdrücklich eingesetzt; eine Diskussion allein über Zentralzahl, Vakuum oder Raumzeitdimension würde den Auftrag verfehlen.

\subsection{Die gemeinsame Mechanik in fünf Schritten}
\begin{enumerate}
\item \textbf{Die Marken legen ein gemeinsames Gerüst fest.} Der P1/P2-Bootstrap, die $3+2$-Ladungszerlegung, der Familienrang und die $D_5+A_3$-Verklebung beschränken dieselbe Architektur. Ihre Auswahlklassen bleiben Teil der Aussage.
\item \textbf{Dieses Gerüst trägt mehrere präzise Auslesungen.} Die volle Materiespur zählt $48$ Komponenten; die Ladungsspur ergibt mit dem Higgs $41/10$. Die rechte schwache Singulettauswahl liefert die drei Flavor-Cusps $1,2/3,1/3$. Das sind verschiedene Operationen auf demselben Träger, keine drei unabhängig angepassten Modelle.
\item \textbf{Wege erzeugen Antworten.} Zwei ursprüngliche Fuchs-Halbwege realisieren den Sechsertransport. Seine gemeinsame Dreikanal-Determinante besitzt den ursprünglichen stationären Flavorpol. Die heute ausgeschriebene Spurabbildung erklärt exakt, wie dieser Transport die vollständigen Farbmultiplets benutzt und dennoch jeden Kanaltyp einmal zählt.
\item \textbf{Ausblenden muss Rückwirkung erhalten.} Die Schur-Reduktion des Sechserwegs lässt den vollständigen Umlauf im verbliebenen Kanal weiterwirken. Derselbe mathematische Grundsatz verbindet an anderer Stelle schwere Neutrinomoden, leichte Masse, Determinantenbeitrag und Gedächtnisantwort. Der Universalraum ist hier die Darstellung, die solche späteren Antworten erhält.
\item \textbf{Die volle physische Ausführung muss gemeinsam folgen.} Die markierte E8-Quelle berechnet bereits gemischte Antworten. Noch muss ihr ursprünglicher Operatoranschluss erklären, warum genau die vorhandenen Flavor-, Naht- und Geometrievariationen auftreten und wie sie lokale Raumzeit, Zustand und Messgrößen erzeugen.
\end{enumerate}

\picturetext{Das Instrument hat mehrere Anzeigen. Eine zählt alle Saiten; eine gewichtet sie nach ihrer Ladung; eine fasst gleichartige Saiten zusammen. Diese Anzeigen dürfen verschiedene Zahlen liefern und dennoch zu einem Instrument gehören. Aber die innere Mechanik muss erklären, welche Bewegung welche Anzeige verändert. Genau diese gemeinsame Mechanik ist der noch ausstehende physische Schluss.}

